\documentclass[a4paper, 12pt]{article}
\usepackage[utf8]{inputenc}
\usepackage[left=3cm, right=3cm]{geometry}
\usepackage[portuguese]{babel}
\usepackage{setspace}
\usepackage{amsmath,amssymb}
\usepackage{multirow}
\usepackage{float}
\usepackage{color}
\usepackage[table,xcdraw]{xcolor}
\usepackage{graphicx}
\usepackage{etoolbox}
\usepackage[hyphens]{url}
\usepackage{hyperref}
\usepackage{pdfpages}
\usepackage{indentfirst}
\usepackage{verbatim}
\usepackage{subfigure}
\setlength{\marginparwidth}{2cm} % ou um valor maior
\usepackage{todonotes}
\usepackage{soul} % no preâmbulo
\usepackage[natbib=true,style=abnt]{biblatex}

\addbibresource{tcc_referencias.bib} % Substitua pelo nome do seu arquivo .bib

% Remove o título da seção de referências
\patchcmd{\thebibliography}{\section*{\refname}}{}{}{}

\newcommand{\citeg}[1]{\textsuperscript{[#1]}}

\hypersetup{
    colorlinks=true,
    linkcolor=black,
    urlcolor=blue,
    citecolor=black
}

\urlstyle{same}

\begin{document}

\onehalfspacing

\thispagestyle{empty}

\begin{center}
{\Huge Trabalho de Conclusão de Curso}
\end{center}

\vspace{3.5cm}

\begin{center}
    \line(1,0){15 cm}
    
    \vspace{1cm}
    
    {\large {\bf Otimização do corte em confecções de roupas}}
    
    \vspace{0.8cm}
    
    \line(1,0){15 cm}

\vspace{3cm}

{\large \textbf{Danillo Mendes Santiago}}

\vspace{3cm}

\noindent
\textbf{Orientadora: Franklina M. B. Toledo}
\end{center}

\vspace{1cm}

\begin{center}   
    \noindent
    \textit{Instituto de Ciências Matemáticas e de Computação}
    
    \noindent
    \textit{Universidade de São Paulo}
    
    \vspace{0.5cm}
    
    \textit{Dezembro/2025}
\end{center}

\newpage
\tableofcontents
\newpage

\section{Introdução}
\label{sec:introducao}

\subsection{Contexto}
\label{sub:contexto}
\todo{ligar as coisas}
A confecção de tecidos é uma atividade humana, presente desde tempos imemoriais e essencial para a produção de vestimentas e utensílios.

Ao longo da história, o desafio de cortar e lidar com materiais permaneceu presente desde a antiguidade sendo uma prática cotidiana de tapeceiros, costureiros e pequenos produtores que lidavam com a disposição e aproveitamento de tecidos em escala artesanal. Com o passar dos séculos, essa atividade foi gradualmente adaptada e transformada em produção industrial, dando origem a um setor de grande relevância econômica. \todo{hono}

Atualmente o processo produtivo pode ser resumido em quatro\todo{ref} etapas: concepção e modelagem, corte das partes que compõem as roupas, costura e acabamento. 
\todo{hono}

A etapa de corte exige a elaboração de diversos planos (arranjos de peças) para atender à demanda de forma eficiente. O planejamento envolve a escolha das peças a serem cortadas em conjunto, o infesto (camadas de tecido) na mesa de corte, o descanso do tecido, o posicionamento das peças no tecido e, finalmente, a execução do corte.\todo{hono}

\subsection{Pergunta Central}

O desafio central deste trabalho, intitulado \textit{Otimização do corte em confecções de roupas}, é minimizar a matéria-prima utilizada para atender a toda a demanda. Para isso, o estudo é voltado ao problema da disposição dos produtos finais no tecido, primeiramente verificando a faixa mínima de tecido, tal que caiba um produto final completo, e posteriormente a escolha da disposição das faixas no tecido final (bin). Com isso é proposta uma solução para a pergunta: \textbf{como dispor e cortas os itens e produtos finais, desperdiçando pouco tecido?}
\todo{hono}

\subsection{Respostas atuais}
Estudos, soluções, Hono, Toledo. exemplos albano, blazewics e marques

\subsection{Objetivo}

Propor métodos heuristicos, para estudo do problema, visando achar boas soluções iniciais, cumprindo as restrições de produção industrial.

\subsection{Relevância}

\todo{continuidade}
Por exemplo, em 2021,\todo{ref} a indústria têxtil faturou aproximadamente R\$ 190 bilhões, representando cerca de 6\% do faturamento industrial e empregando diretamente 1,34 milhões de pessoas. Em 2023, registrava um faturamento de cerca de R\$ 203{,}9 bilhões, evidenciando assim a importância estratégica do setor, diante desse cenário, torna-se natural estudar o problema do corte de tecido pelo viés da Pesquisa Operacional

% \subsection{Estrutura}\todo{se sobrar tempo}

\section{Desenvolvimento}
\label{sec:defprob}

\subsection{Definição do Problema} \todo{hono}

O problema abordado neste estudo é diretamente influenciado pelas \textbf{restrições industriais} do setor. Os produtos finais são compostos por diferentes itens, como mangas, frente e costas de uma camiseta. Uma característica fundamental é que \textbf{os itens de um mesmo produto final devem ser cortados em um único bin} (rolo de tecido ou objeto de corte). Esta restrição é crucial por dois motivos:
\begin{enumerate}
    \item \textbf{Qualidade e padronização:} Garante uniformidade nas roupas, \textbf{evitando diferenças de tonalidade do tecido} entre peças de um mesmo produto.
    \item \textbf{Fluxo de produção:} Permite que o \textbf{processo de costura seja iniciado imediatamente} após o corte do bin. Sem isso, a linha de produção pode sofrer atrasos, pois a costura só começa quando todas as peças de uma roupa estão disponíveis.
\end{enumerate}

Dessa forma, o estudo concentra-se na otimização do corte de peças irregulares em bins, atendendo à restrição de agrupamento, caracterizando-se como um problema de \textbf{empacotamento bidimensional de peças irregulares em um bin}.



\subsection{Classificação do Problema}
\label{sub:classificacao}

NP-Hard

Combinatória

C&P

ISPP

\subsection{Histórico de abordagens e solução?? Revisão da Literatura}
\label{sub:historico}

\section{Metodologia}
\label{sec:metodologia}


\subsection{Modelo Proposto}
\label{sub:modelo}

Dadas as restrições

Explicação Geral, baseado no dotted board, usando meta heuristicas

- BRKGA

- BL

Explicação detalhada


\subsection{BRKGA}
% \label{sub:modelo-base}

O Biased Random-Key Genetic Algorithm (BRKGA) é uma meta-heurística \todo{explicar} evolutiva usada para resolver problemas de otimização combinatória.

- Indivíduo e População, mutação, elite, herança

- Processo evolutivo

- geral



\subsection{BL}
\label{sub:modelo-base}

O Bottom Left (BL) é uma regra de posicionamento construtiva, utilizada em C&P. É uma heuristica, com escolha gananciosa.

- ifp

- nfp

\section{Testes Computacionais}
\label{sec:testes}

\subsection{Geração das Instâncias}
\label{sub:instancias}

\subsection{Análise dos Resultados}
% \label{sub:resultados}

\subsection{Discussão}
% \label{sub:resultados}

- Comparação com Hono e/ou outros

\section{Conclusão}
\label{sec:conclusoes}

\subsection{Síntese}
% \label{sub:resultados}

\subsection{Resposta}
% \label{sub:resultados}

\subsection{Limitações}
\label{sub:limitacoes}

- fica lento com instancias grandes

- fica lento com malha muito fina

% Inserir as referências bibliográficas
\newpage
\printbibliography[title={Referências Bibliográficas}]

\section{Apêndices}
\label{sec:agradecimentos}

\end{document}